\documentclass[a4paper,11pt,fleqn]{article}%fleqn pour aligner les equations à gauche
\usepackage{../../Styles/Cours}


\newcommand{\chnum}{Chapitre 1}
\newcommand{\chtitle}{Identification d'espèces chimiques}
\newcommand{\dtype}{TP}

\begin{document}
\maketitle

\emph{Comment identifier des espèces chimiques ?}
\begin{figure}[h!]
	\begin{minipage}{.48\linewidth}
	\caption{Une situation à éclaircir}
	Lord d'un rangement, le préparateur de laboratoire a retrouvé une série de 4 flacons sans étiquette contenant tous un liquide de nature différente. Comment identifier chaque échantillon ?
	\end{minipage}\hspace{.04\linewidth}
	\begin{minipage}{.48\linewidth}
	\caption{Liste du matériel disponible}
	\begin{itemize*}
		\item Une balance à \SI{0,01}{g} près
		\item 6 béchers
		\item Une éprouvette graduée de \SI{25}{mL}
		\item Une pipette graduée de \SI{10}{mL}
	\end{itemize*}
	\end{minipage}
	
	\vspace{2em}
	
	\begin{minipage}{.48\linewidth}
	\caption{Quelques données}
	\begin{tabular}{|m{.3\linewidth} |m{.2\linewidth} |m{.3\linewidth} |}
	\hline
	Espèce chimique & Masse volumique (\unit{\gram\per\centi\meter\cubed}) & Aspect à température ambiante \\
	\hline
	Eau & \num{1,00} & Liquide incolore\\
	\hline
	Eau salée & \num{1,02} & Liquide incolore\\
	\hline
	Éthanol & \num{0,789} & Liquide incolore\\
	\hline
	Acétone & \num{0,784} & Liquide incolore\\
	\hline
	Glycérol & \num{1,26} & Liquide incolore\\
	\hline
	\end{tabular}
	\end{minipage}\hspace{.04\linewidth}
	\begin{minipage}{.48\linewidth}
	\caption{A propos de la masse volmumique}
	La masse volumique $\rho$ d'un échantillon de matière est une grandeur égale au quotidien de sa masse $m$ par le volume $V$ qu'il occupe.
	Elle est définie par la relation : $$ \rho = \dfrac{m}{V}$$
	avec :
	\begin{itemize**}
		\item la masse $m$ en gramme (\unit{\gram})
		\item le volume $V$ en centimètre cube (\unit{\centi\meter\cubed})
		\item la masse volumique $\rho$ en gramme par centimètre cube (\unit{\gram\per\centi\meter\cubed})
	\end{itemize**}
	La masse volumique est une grandeur qui caractérise une espèce chimique : elle dépend de l'état de l'espèce chimique (gaz, liquide, solide) et de la température.
	\end{minipage}
\end{figure}

\textbf{\textsc{Consigne : En vous aidant des documents fournis et de vos connaissances, répondre aux questions suivantes.}}

\section*{Première Partie :}
\begin{enumerate}
	\item Identifier les grandeurs physiques qu’il faut mesurer pour déterminer la masse volumique d’un échantillon.
	\item Proposer un protocole expérimental permettant de déterminer la masse volumique d'un des  échantillons. Ce protocole doit comporter :
	\begin{itemize*}
		\item les étapes pour réaliser les mesures
		\item les calculs qui devront être effectués
		\item les mesures de sécurité à prendre compte-tenu des pictogrammes de danger indiqués.
		\item la liste du matériel à utiliser
	\end{itemize*}
\end{enumerate}
\begin{center}
		\textbf{Après avoir rédigé votre protocole, appeler le professeur pour le valider.}
	\end{center}

\begin{enumerate}[resume]
	\item Réaliser les mesures et identifier les liquidesà disposition.
	\item La mesure de la masse volumique d’un échantillon permet-elle de déterminer avec certitude sa nature à chaque fois ? Pourquoi ?
\end{enumerate}

\section*{Deuxième Partie :}

\begin{figure}[h!]
	\begin{tabular}{|m{.45\linewidth}|m{.45\linewidth}|}
	\hline
	\caption{Comment prouver la présence d'ions chlorure ?} & \caption{Comment prouver la présence d'éthanol ?}\\
	Les ions chlorures (\ce{Cl-}) sont présents dans le sel commun (appelé aussi chlorure de sodium). Lorsqu’ils sont présents en solution, on peut les identifier en réalisant un test avec une solution de nitrate d’argent. Le test est positif si on observe un précipité blanc quand on verse quelques gouttes de nitrate d’argent dans un tube à essai contenant une solution d’ions chlorures. & On peut prouver la présence d'éthanol grâce à une solution acidifiée de permanganate de potassium. 
Pour réaliser le test, il faut placer dans un tube à essai 2 mL du liquide à tester et y ajouter quelques gouttes de la solution contenant le permanganate de potassium.
Le test est positif si après avoir bouché le tube et l'avoir placé dans un bain marie chaud on observe une décoloration du mélange.\\
\hline	
	\end{tabular}
\end{figure}

\begin{enumerate}[resume]
	\item Quel test, vu au collège, permet d’identifier la présence d’eau ? Proposer un protocole expérimental pour le mettre en \oe uvre.
	\item Mettre en \oe uvre le protocole proposé avec chaque  échantillon et noter vos observations.
    \item Les résultats obtenus permettent-ils de confirmer les conclusions de la première partie ?
    \item Mettre en œuvre le test décrit dans le document 5 avec chaque  échantillon et notez vos observations.
    \item Quelles conclusions peut-on en tirer pour compléter et/ou confirmer l’identification des quatre échantillons ?
    \item Mettre en œuvre le test décrit dans le document 6 avec chaque  échantillon et notez vos observations.
    \item Quelles conclusions peut-on en tirer pour compléter/confirmer l’identification des trois échantillons ?
    \item Parmi ces trois échantillons, lesquels sont des corps purs ? Justifier la réponse.
    \item Parmi ces trois échantillons, lesquels correspondent à un mélange homogène ? Justifier la réponse.
\end{enumerate}
\end{document}